\subsection{Limitations of the Evidence Base}

The strength of these conclusions is limited by evidence quality across four dimensions: sample size constraints, validation setting bias, inconsistent metrics reporting, and insufficient monitoring durations.

\subsubsection{Sample Size Constraints}

The evidence base is constrained by small sample sizes in multiple studies (see Results). Five studies have sample sizes below 20 participants, including single-patient case studies and 3-patient evaluations. Small samples limit the reliability of performance estimates and reduce confidence that results will generalize to broader epilepsy populations. The forecasting literature is particularly affected, with studies limited to 6-11 patients due to invasive EEG confirmation requirements.

\subsubsection{Validation Setting Bias}

Eight of 13 studies were conducted in hospital or EMU settings, creating a systematic bias that may underestimate real-world FPR. Hospital environments lack the diversity of daily activities that generate false alarms in home settings. The substantial FPR difference between EMU validation (0.0054/h) and home validation (0.165/h) demonstrates that hospital-based results may not translate to real-world use.

\subsubsection{Reporting Inconsistency}

Inconsistent metrics reporting hinders cross-study comparison (see Results). Forecasting studies routinely report patient-level success rates, enabling assessment of which patients benefit from the technology. Detection studies rarely provide this granularity, with only two of nine studies reporting individual patient outcomes. This reporting gap prevents assessment of whether detection algorithms work reliably for all patients or only a subset. High aggregate sensitivity may mask substantial variability across patients that remains invisible without patient-level reporting.

\subsubsection{Monitoring Duration and Cold-Start Problems}

Patient-specific approaches require weeks to months of individual data before becoming useful, creating a cold-start problem for newly diagnosed patients. This requirement conflicts with the clinical need for immediate protection after diagnosis. While global models enable immediate deployment, they achieve substantially lower patient success rates. The trade-off between immediate deployment and individualized performance remains unresolved without evidence from hybrid approaches.
